\section{Introduction}
\label{sec:introduction}

This report consolidates the first experimental phase of the IBM Quantum
Credits project devoted to benchmarking Grover--Rudolph state preparation for
amplitude-estimation-based numerical integration on IBM quantum hardware.

The practical question addressed in this phase is not whether asymptotic
quadratic speedups are visible on present-day devices. Rather, the objective is
to establish a reproducible methodology for deciding which amplitude-estimation
circuits are credible candidates for hardware execution, how compilation choices
affect their observed accuracy, and how small structured benchmarks can be used
to expose failure modes before larger Grover--Rudolph instances are attempted.

The campaign follows the feedback received with the IBM Quantum Credits award.
In particular, it combines a QAE-oriented benchmark perspective
\cite{maronese2024qae} with Benchpress-style attention to circuit construction,
transpilation, depth, and two-qubit-gate metrics \cite{nation2025benchpress}.
It also builds on the project's own theoretical foundation: a rigorous
Grover--Rudolph state-preparation analysis \cite{falco2026grproof}, an
angle-structure characterization of quantum numerical integration
\cite{chinesta2026encoding}, an intrinsic gate/compilation viewpoint based on
Lie group embeddings \cite{falco2026liegates}, and the broader algebraic
probability framework in which quantum computation is naturally formulated
\cite{falco2026algebraicprobability}.
The initial hardware target is IBM's Heron backend \texttt{ibm\_fez}. The first
accepted result set demonstrates a complete evidence chain:
\[
  g_0\ \text{calibration}
  \quad\longrightarrow\quad
  g_1\ \text{affine stress test}
  \quad\longrightarrow\quad
  g_2\ \text{quadratic stress test}.
\]
All accepted hardware estimates are produced with a reproducible audit trail,
including logical circuits, transpiled QASM, backend metadata, job identifiers,
counts, and post-processing tables.
